% ---- score2 standard preamble (from: python3 score2_tools.py preamble) ----
\documentclass[10pt]{extarticle}
\usepackage[a4paper,margin=1.5cm]{geometry}
\usepackage[T1]{fontenc}
\usepackage{amsmath,amssymb}
\usepackage{graphicx}
\usepackage[export]{adjustbox}
\usepackage{multicol}
\usepackage{enumitem}
\usepackage{booktabs,array,tabularx}
\setlength{\columnsep}{0.7cm}
\setlength{\parindent}{0pt}
\setlength{\parskip}{3pt}
\pagestyle{empty}
\makeatletter
% figure registry (figdefs.tex is generated by the extract step)
\newcommand{\qfigdef}[4]{\@namedef{qf@file@#1}{#2}\@namedef{qf@w@#1}{#3}\@namedef{qf@wide@#1}{#4}}
\newcommand{\qf@check}[1]{\@ifundefined{qf@file@#1}{\PackageError{score2}{Unknown figure id #1}{}}{}}
\newcommand{\qf@img}[4][]{\includegraphics[width=#2\dimexpr\@nameuse{qf@w@#4}pt\relax,max width=#3,max height=0.45\textheight,#1]{\@nameuse{qf@file@#4}}}
\newif\ifqf@full
\newif\ifqf@boxed
\newsavebox{\qf@box}
\newcommand{\qf@open}{\global\qf@fulltrue\par\global\qf@boxedtrue
  \setbox\qf@box\vbox\bgroup\hsize\textwidth\linewidth\textwidth}
\newcommand{\qf@emit}[1]{\par\egroup\global\qf@boxedfalse
  \ifdim\dimexpr\ht\qf@box+\dp\qf@box\relax<#1\textheight\noindent\usebox{\qf@box}\par\else\unvbox\qf@box\fi}
\newcommand{\qfullsplit}{\ifqf@boxed\qf@emit{0.85}\fi}
\newcommand{\qf@close}{\ifqf@boxed\qf@emit{0.45}\fi\global\qf@fullfalse}
\newcommand{\qfullbegin}{\end{multicols}\qf@open}
\newcommand{\qfullend}{\qf@close\begin{multicols}{2}}
\newcommand{\qfullstart}{\qf@open}
\newcommand{\qfulldone}{\qf@close}
\newcommand{\qfig}[2][1]{\qf@check{#2}%
  \ifqf@full\qf@blk{#1}{#2}\else
  \ifnum\@nameuse{qf@wide@#2}=1 \expandafter\@firstoftwo\else\expandafter\@secondoftwo\fi
  {\end{multicols}\begin{center}\qf@img{#1}{\textwidth}{#2}\end{center}\begin{multicols}{2}}%
  {\qf@blk{#1}{#2}}\fi}
\newcommand{\qf@blk}[2]{\par\smallskip{\centering\qf@img{#1}{\linewidth}{#2}\par}\smallskip}
\newcommand{\qfigi}[2][1]{\qf@check{#2}\qf@img[valign=c]{#1}{\linewidth}{#2}}
\makeatother
\newcommand{\papersec}[1]{\par\bigskip{\centering\Large\bfseries PAPER-#1\par}\medskip}
\newcommand{\partsec}[2]{\par\medskip\noindent{\large\bfseries PART-#1: #2}\par\smallskip}
\newcommand{\sect}[1]{\par\smallskip\noindent\textbf{#1}\par}
\newcommand{\qn}[1]{\par\smallskip\noindent\textbf{#1.}~}
\newcommand{\paraq}[2]{\par\smallskip\noindent\textbf{Paragraph for Questions #1 and #2}\par}
\newenvironment{qoptions}{\begin{itemize}[leftmargin=1.8em,itemsep=1pt,topsep=2pt,label={}]}{\end{itemize}}
\newcommand{\opt}[2]{\item[(#1)] #2}
% figure registry (generated by score2_tools.py finalize)
\qfigdef{p1_phy_s1_q04_1}{Score-II - Test 06 (05 Jul 2022)-img/p1_phy_s1_q04_1.png}{212.5}{0}
\qfigdef{p1_phy_s1_q05_1}{Score-II - Test 06 (05 Jul 2022)-img/p1_phy_s1_q05_1.png}{119.6}{0}
\qfigdef{p1_phy_s2_q08_1}{Score-II - Test 06 (05 Jul 2022)-img/p1_phy_s2_q08_1.png}{193.2}{0}
\qfigdef{p1_phy_s2_q11_1}{Score-II - Test 06 (05 Jul 2022)-img/p1_phy_s2_q11_1.png}{182.0}{0}
\qfigdef{p1_phy_s3_q01_1}{Score-II - Test 06 (05 Jul 2022)-img/p1_phy_s3_q01_1.png}{90.6}{0}
\qfigdef{p1_phy_s3_q05_1}{Score-II - Test 06 (05 Jul 2022)-img/p1_phy_s3_q05_1.png}{240.1}{0}
\qfigdef{p1_chem_s1_q02_1}{Score-II - Test 06 (05 Jul 2022)-img/p1_chem_s1_q02_1.png}{145.4}{0}
\qfigdef{p1_chem_s1_q05_1}{Score-II - Test 06 (05 Jul 2022)-img/p1_chem_s1_q05_1.png}{219.0}{0}
\qfigdef{p1_chem_s1_q05_2}{Score-II - Test 06 (05 Jul 2022)-img/p1_chem_s1_q05_2.png}{222.7}{0}
\qfigdef{p1_chem_s1_q06_1}{Score-II - Test 06 (05 Jul 2022)-img/p1_chem_s1_q06_1.png}{161.1}{0}
\qfigdef{p1_chem_s2_q08_1}{Score-II - Test 06 (05 Jul 2022)-img/p1_chem_s2_q08_1.png}{128.9}{0}
\qfigdef{p1_chem_s2_q09_A}{Score-II - Test 06 (05 Jul 2022)-img/p1_chem_s2_q09_A.png}{95.4}{0}
\qfigdef{p1_chem_s2_q09_B}{Score-II - Test 06 (05 Jul 2022)-img/p1_chem_s2_q09_B.png}{95.0}{0}
\qfigdef{p1_chem_s2_q09_C}{Score-II - Test 06 (05 Jul 2022)-img/p1_chem_s2_q09_C.png}{97.5}{0}
\qfigdef{p1_chem_s2_q09_D}{Score-II - Test 06 (05 Jul 2022)-img/p1_chem_s2_q09_D.png}{97.5}{0}
\qfigdef{p1_chem_s2_q10_1}{Score-II - Test 06 (05 Jul 2022)-img/p1_chem_s2_q10_1.png}{92.5}{0}
\qfigdef{p1_chem_s2_q11_1}{Score-II - Test 06 (05 Jul 2022)-img/p1_chem_s2_q11_1.png}{163.5}{0}
\qfigdef{p1_chem_s2_q11_A}{Score-II - Test 06 (05 Jul 2022)-img/p1_chem_s2_q11_A.png}{53.6}{0}
\qfigdef{p1_chem_s2_q11_B}{Score-II - Test 06 (05 Jul 2022)-img/p1_chem_s2_q11_B.png}{43.5}{0}
\qfigdef{p1_chem_s2_q11_C}{Score-II - Test 06 (05 Jul 2022)-img/p1_chem_s2_q11_C.png}{51.4}{0}
\qfigdef{p1_chem_s2_q11_D}{Score-II - Test 06 (05 Jul 2022)-img/p1_chem_s2_q11_D.png}{53.6}{0}
\qfigdef{p1_chem_s2_q12_A}{Score-II - Test 06 (05 Jul 2022)-img/p1_chem_s2_q12_A.png}{66.3}{0}
\qfigdef{p1_chem_s2_q12_B}{Score-II - Test 06 (05 Jul 2022)-img/p1_chem_s2_q12_B.png}{71.5}{0}
\qfigdef{p1_chem_s2_q12_C}{Score-II - Test 06 (05 Jul 2022)-img/p1_chem_s2_q12_C.png}{68.9}{0}
\qfigdef{p1_chem_s2_q12_D}{Score-II - Test 06 (05 Jul 2022)-img/p1_chem_s2_q12_D.png}{71.5}{0}
\qfigdef{p2_phy_s1_q01_1}{Score-II - Test 06 (05 Jul 2022)-img/p2_phy_s1_q01_1.png}{87.5}{0}
\qfigdef{p2_phy_s1_q02_1}{Score-II - Test 06 (05 Jul 2022)-img/p2_phy_s1_q02_1.png}{92.6}{0}
\qfigdef{p2_phy_s1_q03_1}{Score-II - Test 06 (05 Jul 2022)-img/p2_phy_s1_q03_1.png}{92.1}{0}
\qfigdef{p2_phy_s1_q04_1}{Score-II - Test 06 (05 Jul 2022)-img/p2_phy_s1_q04_1.png}{128.6}{0}
\qfigdef{p2_phy_s1_q05_1}{Score-II - Test 06 (05 Jul 2022)-img/p2_phy_s1_q05_1.png}{126.6}{0}
\qfigdef{p2_phy_s2_q01_1}{Score-II - Test 06 (05 Jul 2022)-img/p2_phy_s2_q01_1.png}{105.5}{0}
\qfigdef{p2_phy_s2_q02_1}{Score-II - Test 06 (05 Jul 2022)-img/p2_phy_s2_q02_1.png}{130.6}{0}
\qfigdef{p2_phy_s2_q03_1}{Score-II - Test 06 (05 Jul 2022)-img/p2_phy_s2_q03_1.png}{119.0}{0}
\qfigdef{p2_phy_s2_q04_1}{Score-II - Test 06 (05 Jul 2022)-img/p2_phy_s2_q04_1.png}{161.8}{0}
\qfigdef{p2_phy_s2_q06_1}{Score-II - Test 06 (05 Jul 2022)-img/p2_phy_s2_q06_1.png}{123.9}{0}
\qfigdef{p2_phy_s3_q03_1}{Score-II - Test 06 (05 Jul 2022)-img/p2_phy_s3_q03_1.png}{124.6}{0}
\qfigdef{p2_phy_s3_q04_1}{Score-II - Test 06 (05 Jul 2022)-img/p2_phy_s3_q04_1.png}{85.2}{0}
\qfigdef{p2_phy_s3_q05_1}{Score-II - Test 06 (05 Jul 2022)-img/p2_phy_s3_q05_1.png}{104.9}{0}
\qfigdef{p2_chem_s1_q01_1}{Score-II - Test 06 (05 Jul 2022)-img/p2_chem_s1_q01_1.png}{66.3}{0}
\qfigdef{p2_chem_s1_q02_1}{Score-II - Test 06 (05 Jul 2022)-img/p2_chem_s1_q02_1.png}{203.3}{0}
\qfigdef{p2_chem_s1_q02_A}{Score-II - Test 06 (05 Jul 2022)-img/p2_chem_s1_q02_A.png}{53.6}{0}
\qfigdef{p2_chem_s1_q02_B}{Score-II - Test 06 (05 Jul 2022)-img/p2_chem_s1_q02_B.png}{53.6}{0}
\qfigdef{p2_chem_s1_q02_C}{Score-II - Test 06 (05 Jul 2022)-img/p2_chem_s1_q02_C.png}{53.1}{0}
\qfigdef{p2_chem_s1_q02_D}{Score-II - Test 06 (05 Jul 2022)-img/p2_chem_s1_q02_D.png}{53.6}{0}
\qfigdef{p2_chem_s2_q01_1}{Score-II - Test 06 (05 Jul 2022)-img/p2_chem_s2_q01_1.png}{67.1}{0}
\qfigdef{p2_chem_s2_q04_1}{Score-II - Test 06 (05 Jul 2022)-img/p2_chem_s2_q04_1.png}{45.1}{0}
\qfigdef{p2_chem_s3_q01_1}{Score-II - Test 06 (05 Jul 2022)-img/p2_chem_s3_q01_1.png}{270.1}{1}
\qfigdef{p2_chem_s3_q02_1}{Score-II - Test 06 (05 Jul 2022)-img/p2_chem_s3_q02_1.png}{328.2}{1}
\begin{document}
\begin{center}{\Large\bfseries JEE (Main+Advanced) Score-II --- Test 06 (05 Jul 2022)}\end{center}
\begin{multicols}{2}

\papersec{1}
\partsec{1}{PHYSICS}
\sect{Section-I (i): Single Correct}

\qn{1} A capacitor of capacitance $C_1 = 3.0\ \mu F$ withstands the maximum voltage $V_1 = 4.0$ kV, while a capacitor of capacitance $C_2 = 6.0\ \mu F$ the maximum voltage $V_2 = 3.0$ kV. If they are connected in series what maximum voltage can the system withstand?
\begin{qoptions}
\opt{A}{3.0 kV}
\opt{B}{4.0 kV}
\opt{C}{6.0 kV}
\opt{D}{7.0 kV}
\end{qoptions}

\qn{2} A series circuit is made up of two cells of emf 1.5 and 3V, respectively, and two coils each of resistance $10\Omega$, arranged in the order cell, coil, cell, coil. The centre points of the two coils are joined by a sensitive galvanometer which shows no deflection. If the cell of 1.5V has an internal resistance of $5\Omega$, calculate the internal resistance of the other cell.
\begin{qoptions}
\opt{A}{$5\Omega$}
\opt{B}{$10\Omega$}
\opt{C}{$15\Omega$}
\opt{D}{$20\Omega$}
\end{qoptions}

\qn{3} A pair of circular coils each having 50 turns of radius 50 cm are separated by 50 cm. A current of 10A passes through the coils which are connected in series. Midway between the coils, a coaxial flat metal disc of radius 10cm, is rotating at 1000 rpm. What is the approximate emf generated between the centre and the rim of the disc ($\mu_0 = 4\pi \times 10^{-7}$ H/m)?
\begin{qoptions}
\opt{A}{1.0 mv}
\opt{B}{0.5 mv}
\opt{C}{2.0 mv}
\opt{D}{4.0 mv}
\end{qoptions}

\qn{4} Water is moving through a pipe. The velocity profile at some section is shown fig. and is given mathematically as $v = \left(\dfrac{\beta}{4\mu}\right)\left(\dfrac{d^2}{4} - r^2\right)$, where $v$ = velocity of water at any position $r$, $\beta$ = a constant, $\mu$ = viscosity of water, $d$ = pipe diameter, and $r$ = radial distance from centerline. If the given profile persists a distance $L$ along the pipe, what drag is induced on the pipe by the water in the direction of flow over this distance?
\qfig{p1_phy_s1_q04_1}
\begin{qoptions}
\opt{A}{$\dfrac{\beta d^2 \pi L}{8}$}
\opt{B}{$\dfrac{\beta d^2 \pi L}{4}$}
\opt{C}{$\beta d^2 \pi L$}
\opt{D}{Zero}
\end{qoptions}

\qn{5} Two blocks of masses 'm' and 'M' are connected by a chord passing around a frictionless pulley which is attached to a rotating frame, which rotates about a vertical axis with an angular velocity $\omega$.
\qfig{p1_phy_s1_q05_1}
If the coefficient of friction between the two masses and the surface be $\mu_1$ and $\mu_2$ respectively, determine the value of $\omega$, at which the block starts sliding radially (M > m).
\begin{qoptions}
\opt{A}{$\left[\dfrac{g(\mu_1 m + \mu_2 M)}{Mr_2 + mr_1}\right]^{1/2}$}
\opt{B}{$\left[\dfrac{g(\mu_1 m + \mu_2 M)}{Mr_2 - mr_1}\right]^{1/2}$}
\opt{C}{$\left[\dfrac{g(\mu_1 m - \mu_2 M)}{(Mr_2 - mr_1)}\right]^{1/2}$}
\opt{D}{$\left[\dfrac{g(\mu_1 m - \mu_2 M)}{(Mr_2 + mr_1)}\right]^{1/2}$}
\end{qoptions}

\qn{6} A smooth sphere of mass m = 1 kg, moving with a velocity 'u', collides obliquely on a smooth sphere of mass M = 2 kg at rest, such that, its direction of velocity, before impact makes an angle $\theta$ to the line of impact. At what value of $\theta$, the sphere of mass m will be deflected through a right angle from its initial direction, for elastic collision.
\begin{qoptions}
\opt{A}{$30^\circ$}
\opt{B}{$37^\circ$}
\opt{C}{$45^\circ$}
\opt{D}{$60^\circ$}
\end{qoptions}

\sect{Section-I (ii): One or More Correct}

\qn{7} In a block of cobalt/iron alloy, it is suspected that the cobalt (Z = 27) is very poorly mixed with the iron (Z = 26). Given that the ionization energy of hydrogen is 13.6 eV then the binding energies and their wavelength of the K shell (called as absorption edges) of the constituents of the alloys are:
\begin{qoptions}
\opt{A}{$E_K$ of cobalt = 9.2 Kev}
\opt{B}{$E_K$ of iron = 8.5 Kev}
\opt{C}{$\lambda_K$ of cobalt = 1.35\AA}
\opt{D}{$\lambda_K$ of iron = 1.46 \AA}
\end{qoptions}

\qn{8} Early in the morning, on a sunny day, the heat of the sun produces a thin layer of warm air above the surface of a long straight road causes the formation of mirage. Consider a possible light ray path such as that illustrated in figure. This connects an eye--level point on the tree with an observer of height h = 2m. If the layer of hot air has refractive index $n_2 = 1.00020$, while the cold air has refractive index $n_1 = 1.00030$. Then
\qfig{p1_phy_s2_q08_1}
\begin{qoptions}
\opt{A}{optical path length along ABCD $= n_2 x + n_1\sqrt{(d-x)^2 + 4h^2}$}
\opt{B}{the actual distance that the ray travels in the layer of hot air is 217.2 m when d = 500 m}
\opt{C}{the actual distance that the ray travels in the layer of hot air is 282.8 m when d = 500 m}
\opt{D}{at 100 m from tree mirage disappears.}
\end{qoptions}

\qn{9} A metal surface is illuminated with light of different wavelengths and the corresponding stopping potential of the photoelectrons V, are found to be as follows:
\begin{center}
\begin{tabular}{|c|c|c|c|c|c|c|}
\hline
$\lambda(\text{\AA})$ & 3,660 & 4,050 & 4,360 & 4,920 & 5,460 & 5,790\\
\hline
V(V) & 1.48 & 1.15 & 0.93 & 0.62 & 0.36 & 0.24\\
\hline
\end{tabular}
\end{center}
From given data
\begin{qoptions}
\opt{A}{Planck's constant $h = 4.13\times 10^{-15}$ ev--s}
\opt{B}{Work function of metal = 1.9 ev}
\opt{C}{Threshold frequency $= 4.6\times 10^{15}$ Hz}
\opt{D}{Threshold wavelength $= 6.6\times 10^{-7}$ m}
\end{qoptions}

\qn{10} Choose the correct option(s). Which of the following will change the time period as they are taken to Mars?
\begin{qoptions}
\opt{A}{A simple pendulum}
\opt{B}{A compound pendulum}
\opt{C}{LC circuit}
\opt{D}{A torsional pendulum}
\end{qoptions}

\qn{11} A uniform electric field is established between the plates of a parallel plate capacitor by holding one plate at ground and the other at a positive potential V as shown. A uniform magnetic field B is established perpendicular to the electric field as shown. If a charge $-q$ is released from rest from the lower plate at O. Then (Neglect gravity)
\qfig{p1_phy_s2_q11_1}
\begin{qoptions}
\opt{A}{$V_x = \dfrac{V/d}{B}\left\{1 - \cos\left(\dfrac{qB}{m}t\right)\right\}$}
\opt{B}{$V_y = \dfrac{V/d}{B}\sin\left(\dfrac{qBt}{m}\right)$}
\opt{C}{$V_x = \left(\dfrac{qBy}{m}\right)$}
\opt{D}{$V_y = \sqrt{\dfrac{2qV}{md}y - \dfrac{q^2B^2}{m^2}y^2}$}
\end{qoptions}

\qn{12} In the Kundt's tube experiment, the length of the steel rod which is stroked is 120 cm long and the distance between heaps of cork dust is 8 cm when the rod is caused to vibrate longitudinally in air. If the ends of the tube are sealed and the air replaced by a gas and the experiment repeated, the distance between heaps is observed to be 10 cm. [Assume velocity of sound in air = 340 m/s]. Choose the correct options.
\begin{qoptions}
\opt{A}{Velocity of sound in gas = 425 m/s}
\opt{B}{Velocity of sound in rod = 2550 m/s}
\opt{C}{Frequency of rod = 2125 Hz}
\opt{D}{Velocity of sound in rod is 12 times velocity in gas}
\end{qoptions}

\sect{Section-II: Numerical Value}

\qn{1} A non--conducting ring R having uniformly distributed charge Q starts rotating about x--x' axis passing through diameter with an angular acceleration $\alpha$, as shown in the figure. Another small conducting ring having radius $a (a \ll R)$ is kept fixed at the centre of bigger ring in such a way that axis xx' is passing through its centre and perpendicular to its plane. If the resistance of small ring is $r = 1\Omega$, find the induced current in it in ampere. (Given $Q = \dfrac{16\times 10^2}{\mu_0}$ C, $R = 1$m, $a = 0.1$ m, $\alpha = 8$ rad s$^{-2}$)
\qfig{p1_phy_s3_q01_1}

\qn{2} A beam of monochromatic light wavelength $5.82\times 10^{-7}$ m falls normally on a glass wedge with the wedge angle of 20 s of an arc. If the refractive index of glass is 1.5, find the number of interference fringes per centimetre of the wedge length.

\qn{3} 100 millicuries of radon which emits 5.5 MeV $\alpha$--particles are contained in a glass capillary tube 5 cm long with internal and external diameters 2 and 8 mm respectively. Neglecting end effects and assuming that the inside of the tube is uniformly irradiated by the particles which are stopped at the surface. Calculate the temperature difference in milli degree Celsius between the walls of a tube when steady thermal conditions have been reached.\\
Thermal conductivity of glass = 0.025 Cal cm$^{-1}$s$^{-1}$C$^{-1}$\\
Curie = $3.7\times 10^{10}$ disintegrations per second\\
J = 4.2 joule Cal$^{-1}$

\qn{4} A popular formula in the hydraulics literature is the Hazen--Williams formula for volume flow rate (volume/unit-time) in a pipe of diameter D and pressure gradient $dp/dx$ is $61.9 D^{2.63}(dp/dx)^{0.54}$. If the dimensions of the constant 61.9 is $M^aL^bT^c$ then $a+b+c = ?$

\qn{5} A force of 460 N is exerted on lever AB, as shown in fig. End B is connected to a piston which fits into a cylinder having a diameter of 60 mm. What force $F_D$ in KN acts on the larger piston for equilibrium, if the volume between C and D is filled with water?
\qfig{p1_phy_s3_q05_1}

\qn{6} Two points P and Q moves in a straight line AB. The point P starts from A in the direction AB with velocity $u_1$ and acceleration $f_1$ and at the same time Q starts from B in the direction of BA with velocity $u_2$ and acceleration $f_2$. If they pass one another at the middle point of AB and arrive at the other ends of AB with equal speeds, then $(u_1+u_2)(f_1-f_2) = k(f_1 u_2 - f_2 u_1)$. Find the value of k.

\partsec{2}{CHEMISTRY}
\sect{Section-I (i): Single Correct}

\qn{1} 1 mole each of CaC$_2$, Mg$_2$C$_3$, Al$_4$C$_3$ reacts with H$_2$O in separate open flasks. The magnitude of work done during dissolution shows the order:
\begin{qoptions}
\opt{A}{CaC$_2$ = Mg$_2$C$_3$ < Al$_4$C$_3$}
\opt{B}{CaC$_2$ = Mg$_2$C$_3$ = Al$_4$C$_3$}
\opt{C}{Mg$_2$C$_3$ < CaC$_2$ < Al$_4$C$_3$}
\opt{D}{Mg$_2$C$_3$ < Al$_4$C$_3$ < CaC$_2$}
\end{qoptions}

\qn{2} Consider a following setup in which two glucose solution having concentration $C_1$ and $C_2$ are separated by semipermeable membrane as shown below.
\qfig{p1_chem_s1_q02_1}
At what values of external pressures $P_1$ \& $P_2$, osmosis will occur?
\begin{qoptions}
\opt{A}{$P_1 = 0$, $P_2 = (C_2-C_1)RT$}
\opt{B}{$P_1 = C_1RT$, $P_2 = C_2RT$}
\opt{C}{$P_1 = C_2RT$, $P_2 = C_1RT$}
\opt{D}{$P_1 = 2C_1RT$, $P_2 = (C_1+C_2)RT$}
\end{qoptions}

\qn{3} The oxidation of certain metal is found to obey the equation $A^2 = \alpha t + \beta$, where A is the thickness of metal oxide film at time t, $\alpha$ and $\beta$ are constants. The order of this reaction is
\begin{qoptions}
\opt{A}{0}
\opt{B}{1}
\opt{C}{$-1$}
\opt{D}{2}
\end{qoptions}

\qn{4} An ideal gas undergoes adiabatic expansion against constant external pressure. Which of the following is incorrect: ($T_f$ and $T_i$ are final and initial temperatures respectively)
\begin{qoptions}
\opt{A}{$\Delta U + P_{ext}\Delta V = 0$}
\opt{B}{K.E. of gas sample increases}
\opt{C}{$\Delta S_{surr} = 0$}
\opt{D}{$W = \dfrac{nR}{\gamma - 1}(T_f - T_i)$}
\end{qoptions}

\qn{5} Correct statement among the following is:
\qfig{p1_chem_s1_q05_1}
\qfig{p1_chem_s1_q05_2}
\begin{qoptions}
\opt{A}{Product A and B are related as Identical compounds}
\opt{B}{Product A and C are related as Diastereomers}
\opt{C}{Product D and C are related as Identical compounds}
\opt{D}{Product D and A are related as Diastereomers}
\end{qoptions}

\qn{6} For the given reaction:
\qfig{p1_chem_s1_q06_1}
Correct statement among the following is:\\
\relax[Hint. Use two equivalent of Y for preparation of Z]
\begin{qoptions}
\opt{A}{'Z' product can be obtained on treatment of C$_6$H$_6$ with C$_6$H$_5\overset{\oplus}{\mathrm{N}}{}_2\mathrm{Cl}^{\ominus}$}
\opt{B}{On treatment with Sn/HCl compound 'Z' will form an azo compound.}
\opt{C}{'Y' can be used as a solvent during Friedal Craft Alkylation.}
\opt{D}{Y will give negative unsaturation test with Bromine water}
\end{qoptions}

\sect{Section-I (ii): One or More Correct}

\qn{7} Sometimes yellow turbidity appears while passing H$_2$S gas even in the absence of $2^{nd}$ group of basic radicals. Which of the following is/are not correct justification for the observation:
\begin{qoptions}
\opt{A}{Sulphur is present in mixture as impurity}
\opt{B}{IV group radicals are precipitated as sulphides.}
\opt{C}{H$_2$S is oxidised by some acid radicals}
\opt{D}{III$^{rd}$ group radicals are precipitated as yellow coloured hydroxides}
\end{qoptions}

\qn{8} For the given reaction:
\qfig{p1_chem_s2_q08_1}
Which of the following is/are correct statements regarding gas (D)?
\begin{qoptions}
\opt{A}{Hybridisation of gas is sp$^2$}
\opt{B}{It can act as oxidising agent}
\opt{C}{It gives starch-iodate test}
\opt{D}{It can decolourise KMnO$_4$}
\end{qoptions}

\qn{9} \textbf{Statement-1}: Net work done by system or on the system lies between 100 cal to 200 cal.\\
\textbf{Statement-2}: Net heat released or absorbed by the system lies between 100 cal to 200 cal.\\
Which of the following graph(s) for one mol of an ideal mono atomic gas is/are satisfying both the statements ($\log 2 = 0.30$, $R = \dfrac{2\ \text{cal}}{\text{mol K}}$)?
\begin{qoptions}
\opt{A}{\qfigi{p1_chem_s2_q09_A}}
\opt{B}{\qfigi{p1_chem_s2_q09_B}}
\opt{C}{\qfigi{p1_chem_s2_q09_C}}
\opt{D}{\qfigi{p1_chem_s2_q09_D}}
\end{qoptions}

\qn{10} An orbital in a H-atom is associated with the following radial and angular parts of its wave function
\resizebox{\linewidth}{!}{$\psi(r) = \dfrac{1}{9\sqrt{30}}\left(\dfrac{Z}{a_0}\right)^{3/2}\sigma^2 e^{-\sigma/2};\quad \psi(\theta,\phi) = \left(\dfrac{15}{4\pi}\right)^{1/2}\sin\theta\cos\theta\cos\phi.$}
Here, $\sigma = \dfrac{2r}{na_0}$, $a_0 = \dfrac{h^2}{4\pi^2 me^2}$, $\theta$ = colatitude, $\phi$ = azimuth.
\qfig{p1_chem_s2_q10_1}
Select the correct option(s) regarding this orbital:
\begin{qoptions}
\opt{A}{Total energy of electron in it = $-1.51$ eV}
\opt{B}{Orbital angular momentum of electron in this orbital is $\sqrt{6}\,\hbar$}
\opt{C}{ZX plane is not the nodal plane.}
\opt{D}{Orbital can be certainly called $3d_{zx}$ orbital.}
\end{qoptions}

\qn{11} Consider the following reactions.
\qfig{p1_chem_s2_q11_1}
'D' product can be:
\begin{qoptions}
\opt{A}{\qfigi{p1_chem_s2_q11_A}}
\opt{B}{\qfigi{p1_chem_s2_q11_B}}
\opt{C}{\qfigi{p1_chem_s2_q11_C}}
\opt{D}{\qfigi{p1_chem_s2_q11_D}}
\end{qoptions}

\qn{12} Which of the following is/are Resolvable?
\begin{qoptions}
\opt{A}{\qfigi{p1_chem_s2_q12_A}}
\opt{B}{\qfigi{p1_chem_s2_q12_B}}
\opt{C}{\qfigi{p1_chem_s2_q12_C}}
\opt{D}{\qfigi{p1_chem_s2_q12_D}}
\end{qoptions}

\sect{Section-II: Numerical Value}

\qn{1} How many of the following salts are white and evolves colourless gas or vapours with dil. H$_2$SO$_4$.\\
BaCrO$_4$, Hg$_2$CrO$_4$, ZnS, BaSO$_4$, BaS$_2$O$_3$, CH$_3$COOAg, AgNO$_2$.

\qn{2} Three moles of an ideal gas is changed from its initial state (16 atm, 6L) to final state (4 atm, 15L) in such a way that this change can be represented by a straight line in P-V curve. The maximum temperature attained by the gas during the change is 108Y kelvin. Calculate Y. ($R = 1/12$)

\qn{3} Let Z be the atomic number of the element with the maximum number of unpaired electrons in 4p subshell and having spin multiplicity S. Let n and P are the group number and period number respectively in modern periodic table for this element. Find $\dfrac{Z+P+S+n}{8}$.

\qn{4} How many of the following reagents bleach litmus solution and also produce blue black colour in starch iodide solution. O$_3$, NH$_3$, N$_2$, H$_2$SO$_4$, NO, CO$_2$, H$_2$O$_2$, PCl$_3$, Cl$_2$, Cl$_2$O.

\qn{5} From the given reaction\\
\resizebox{0.95\linewidth}{!}{$\mathrm{C_6H_6-C\equiv C-C_6H_5} \xrightarrow[\text{Me}_3\text{COK}]{\text{C}_6\text{H}_5\text{CHBr}_2}$ 'X'}\\
If total number of carbon atom in 'X' is A and degree of unsaturation in 'X' is B then find out $\dfrac{2\times A}{B} =$

\qn{6} Out of the given amino acids.\\
Arginine, Histidine, Phenyl alkanine, Tyrosine and Tryptophan.\\
$P_1$ Amino acids give +ve Pauly's Test. (Test for imidazole group)\\
$P_2$ Amino acids give +ve Sakaguchi reaction/test. (Test for guanidino group)\\
$P_3$ Amino acids give +ve Xanthoprotic Test.\\
$P_4$ Amino acids give +ve Biuret test.\\
$P_5$ Amino acids give +ve Ninhydrin test.\\
$P_6$ Amino acids give Millon's Reaction. (Test for Phenolic group)\\
The sum of $\displaystyle\sum_{i=1}^{6} P_i = X$. Find $\dfrac{X}{4}$.

\partsec{3}{MATHEMATICS}
\sect{Section-I (i): Single Correct}

\qn{1} Let $f(x)$ be a differentiable function satisfying $xf(x) - \displaystyle\int_0^x f(t)\,dt = x + \ln\left(\sqrt{1+x^2}-x\right)$ and $f(0) = \ln 2$. If $g(x) = xf'(x)$, then range of $g(x)$ is
\begin{qoptions}
\opt{A}{$(-\infty,\infty)$}
\opt{B}{$[1,\infty)$}
\opt{C}{$[0,1)$}
\opt{D}{$(-\infty,0]$}
\end{qoptions}

\qn{2} The minimum distance between the curves $y = \tan^{-1}\left(\dfrac{(x^3-7)\sin^{-1}(1-x)}{\cot^{-1}(\cos^{-1}(3-x))}\right)$ and $(x+4)^2 + \left(y+\dfrac{\pi}{4}\right)^2 = 1$ is equal to
\begin{qoptions}
\opt{A}{4}
\opt{B}{5}
\opt{C}{6}
\opt{D}{7}
\end{qoptions}

\qn{3} Let $n \geq 3$ and let $C_1, C_2, \ldots, C_n$ be circles with radii $r_1, r_2, \ldots, r_n$ respectively where $r_1 < r_2 < r_3 \ldots < r_n$. Assume that $C_i$ and $C_{i+1}$ touch externally where $1 \le i \le n-1$. It is also given that the x-axis and the line $3y = 4x+9$ are tangential to each of the circles. Then $r_1, r_2, \ldots, r_n$ are in
\begin{qoptions}
\opt{A}{An arithmetic progression with common difference $\dfrac{3+\sqrt5}{2}$}
\opt{B}{A geometric progression with common ratio $\dfrac{(\sqrt3+5)}{2}$}
\opt{C}{An arithmetic progression with common difference $\dfrac{(\sqrt3+5)}{2}$}
\opt{D}{A geometric progression with common ratio $\dfrac{(3+\sqrt5)}{2}$}
\end{qoptions}

\qn{4} Let $S(x) = \displaystyle\int \dfrac{dx}{e^x+8e^{-x}+4e^{-3x}}$, $R(x) = \displaystyle\int \dfrac{dx}{e^{3x}+8e^x+4e^{-x}}$ and $M(x) = S(x) - 2R(x)$. If $M(x) = \dfrac{1}{2}\tan^{-1}\left(\dfrac{f(x)}{2}\right) + c$, (where c is arbitrary constant) and $M(0) = \dfrac{1}{2}\tan^{-1}\left(\dfrac{3}{2}\right)$, then $f(\log_e 2)$ is equal to
\begin{qoptions}
\opt{A}{1}
\opt{B}{3}
\opt{C}{5}
\opt{D}{7}
\end{qoptions}

\qn{5} If a tangent of slope 6 to the ellipse $\dfrac{x^2}{a^2} + \dfrac{y^2}{b^2} = 1$ is normal to the circle $x^2+y^2+4x+1=0$ then the maximum value of $ab$ is (where $a>0$, $b>0$)
\begin{qoptions}
\opt{A}{10}
\opt{B}{14}
\opt{C}{12}
\opt{D}{8}
\end{qoptions}

\qn{6} If a and b are distinct positive integers and the quadratic equations $(a-1)x^2 - (a^2+2)x + (a^2+2a) = 0$ and $(b-1)x^2 - (b^2+2)x + (b^2+2b) = 0$ have a common root. Then which of the following can be true.
\begin{qoptions}
\opt{A}{$a^2+b^2 = 45$}
\opt{B}{$a = 2b$}
\opt{C}{$b = 3a$}
\opt{D}{$ab = 18$}
\end{qoptions}

\sect{Section-I (ii): One or More Correct}

\qn{7} Let $z_1, z_2, z_3$ be three distinct complex numbers satisfying $|z_1| = |z_2| = |z_3| = 1$. Which of the following is true?
\begin{qoptions}
\opt{A}{if $\arg\left(\dfrac{z_1}{z_2}\right) = \dfrac{\pi}{2}$ then $\arg\left(\dfrac{z-z_1}{z-z_2}\right) > \dfrac{\pi}{4}$ where $|z|>1$}
\opt{B}{$|z_1z_2 + z_2z_3 + z_3z_1| = |z_1+z_2+z_3|$}
\opt{C}{$\mathrm{Im}\left(\dfrac{(z_1+z_2)(z_2+z_3)(z_3+z_1)}{z_1\cdot z_2\cdot z_3}\right) = 0$}
\opt{D}{If $|z_1-z_2| = \sqrt{2}|z_1-z_3| = \sqrt{2}|z_2-z_3|$, then $\mathrm{Re}\left(\dfrac{z_3-z_1}{z_3-z_2}\right) = 0$}
\end{qoptions}

\qn{8} {\sloppy For the equation $2\cos 2x\left(\cos 2x - \cos\left(\dfrac{2014\pi^2}{x}\right)\right) = (\cos 4x) - 1$; which of the following is/are correct?\par}
\begin{qoptions}
\opt{A}{The number of solutions of the equation is 3.}
\opt{B}{The sum of all the positive solutions is $1080\pi$}
\opt{C}{The number of positive solutions is 4}
\opt{D}{The sum of all positive solutions is $1008\pi$}
\end{qoptions}

\qn{9} If x is number of all rational number $\dfrac{m}{n}$ such that\\
(i) $0 < \dfrac{m}{n} < 1$\\
(ii) H.C.F. $\{m,n\} = 1$\\
(iii) $m\cdot n = 30!$\\
then
\begin{qoptions}
\opt{A}{x is 256}
\opt{B}{x is 512}
\opt{C}{Sum of digits of x is 8}
\opt{D}{Sum of digits of x is 13}
\end{qoptions}

\qn{10} Let $\lambda, \alpha, \beta \in \mathbb{R}$, consider the system of linear equations
$$10x + \lambda y = \alpha$$
$$-5x + 3y = 2\beta$$
then which of following statements is(are) correct?
\begin{qoptions}
\opt{A}{if $\lambda \neq -6$ then system has unique solutions for all values of $\alpha$ and $\beta$}
\opt{B}{if $\alpha + 4\beta = 0$ then system has infinite many solutions for $\lambda = -6$}
\opt{C}{if $\lambda = -6$ then system has infinite many solutions of all values of $\alpha$ and $\beta$}
\opt{D}{if $\alpha + 4\beta \neq 0$ then system has no solutions for $\lambda = -6$}
\end{qoptions}

\qn{11} Let $f(x,y) = \sqrt{x^2+y^2} + \sqrt{x^2+y^2-4x+4} + \sqrt{x^2+y^2-4y+4} + \sqrt{x^2+y^2-12x-16y+100}$ $\forall x,y\in\mathbb{R}$ then
\begin{qoptions}
\opt{A}{Minimum value of $f(x,y) = 10 + 2\sqrt{2}$}
\opt{B}{Minimum value of $f(x,y) = 10 - 2\sqrt{2}$}
\opt{C}{Minimum value of $f(x,y)$ occurs for $x = \dfrac{6}{7}$}
\opt{D}{Minimum value of $f(x,y)$ occurs for $y = \dfrac{8}{7}$}
\end{qoptions}

\qn{12} TP and TQ are tangents to parabola $y^2 = 12x$ and normals at P and Q intersect at point R on the parabola. The locus of centre of Circumcircle of $\triangle$TPQ is a parabola whose
\begin{qoptions}
\opt{A}{length of latus-rectum is $\dfrac{3}{4}$}
\opt{B}{focus is $\left(\dfrac{27}{8},0\right)$}
\opt{C}{vertex is $(3,0)$}
\opt{D}{foot of directrix is $\left(\dfrac{21}{8},0\right)$}
\end{qoptions}

\sect{Section-II: Numerical Value}

\qn{1} Let $I_n = \displaystyle\int_0^{3\pi/2} (\ln|\sin x|)\cos(2nx)\,dx$ where $n\in\mathbb{N}$. If $(12I_3 - 16I_2)$ is $k\pi$ then value of k is

\qn{2} Let $f:\mathbb{R}\to\mathbb{R}$ be a twice differentiable function satisfying $f''(x) - 5f'(x) + 6f(x) = 0\ \forall x\in\mathbb{R}$. If $f(0) = 1$, $f'(0) = 0$ then maximum value of $f(x)$ is

\qn{3} Let $p(x)$ be a polynomial of degree 100 such that $p(x) = \dfrac{1}{(x-4)}$ for $x = 5,6,7,\ldots,104,105$ and $p(106) = \dfrac{m}{n}$ (where m, n are coprime numbers) then value of $(n-m)$ is

\qn{4} Non-zero vectors $\vec a, \vec b, \vec c$ satisfy $\vec a\cdot\vec b = 0$, $(\vec b - \vec a)\cdot(\vec b + \vec c) = 0$ and $2|\vec b+\vec c| = |\vec b - \vec a|$. If $\vec a = \mu\vec b + 4\vec c$ then the value of $\mu$ is

\qn{5} Let $p = \displaystyle\sum_{r=1}^{2021}(-1)^{r-1}\dfrac{{}^{2021}C_r}{2^{r+1}-1}$ and $q = \displaystyle\sum_{r=1}^{\infty}\dfrac{(2^r-1)^{2021}}{2^{2022r}}$, then the value of $p+q$ is

\qn{6} If $\triangle ABC$ is right angle triangle (right angle at B) inscribed in the parabola $y^2 = 4x$, the minimum length of intercept cut off by the tangents at A and C to the parabola from y-axis is?

\papersec{2}
\partsec{1}{PHYSICS}
\sect{Section-I: One or More Correct}

\qn{1} Suppose in an insulating medium having dielectric constant $K = 1$ volume density of positive charge varies with y--co-ordinate according to the law $\rho = 4y$. A particle of mass m having positive charge q is placed in the medium at a point $(0,10)$ and projected with velocity $\vec V = 100\hat i$ as shown in figure. Neglecting the gravity and frictional resistance of the medium and assuming electric field strength to be zero at $y=0$. Choose the correct statements.
\qfig{p2_phy_s1_q01_1}
\begin{qoptions}
\opt{A}{The direction of electric field is along positive y--axis.}
\opt{B}{The direction of electric field is along negative y--direction.}
\opt{C}{The magnitude of electric field E in the region as a function of y is $\dfrac{2y^2}{\epsilon_0}$.}
\opt{D}{The magnitude of electric field E in the region as function of y is $\dfrac{8y^2}{\epsilon_0}$.}
\end{qoptions}

\qn{2} A source of internal resistance $4\Omega$ is connected in a circuit as shown in figure. R is a variable resistance. As value of resistance R changes, current drawn from source as well as power dissipated in each individual resistance also change.
\qfig{p2_phy_s1_q02_1}
\begin{qoptions}
\opt{A}{The value of R for which power consumption in it will be maximum is $5\Omega$.}
\opt{B}{The value of R for which power consumption in it will be maximum is $8\Omega$.}
\opt{C}{If resistance found for maximum power (for case A or B) is used, then terminal potential difference across the source be 0.8V}
\opt{D}{If resistance found for maximum power (for case (A) or (B)) is used then terminal potential difference across the source will be 0.6V}
\end{qoptions}

\qn{3} A large parallel plate capacitor with uniform surface charge $\sigma$ on the upper plate and $-\sigma$ on the lower is moving with a constant speed V as shown in figure
\qfig{p2_phy_s1_q03_1}
\begin{qoptions}
\opt{A}{Magnetic field between the plates is $\mu_0\sigma V$.}
\opt{B}{Magnetic field between the plates is $\left(\dfrac{\mu_0\sigma V}{2}\right)$.}
\opt{C}{Magnetic force per unit area on the upper plate is $\dfrac{\mu_0\sigma^2V^2}{2}$ upward.}
\opt{D}{Magnetic force per unit area on the upper plate is $\mu_0\sigma^2V^2$ upward.}
\end{qoptions}

\qn{4} An infinitely long wire lying along Z-axis carries a current I, flowing towards positive Z-direction. There is no other current, consider a circle in x--y plane with centre at $(2,0,0)$ and radius 1 meter. Divide the circle in small segments let $d\vec\ell$ denote the length segment in anticlockwise direction as shown.
\qfig{p2_phy_s1_q04_1}
\begin{qoptions}
\opt{A}{The path integral $\oint \vec B \cdot d\vec\ell$ of the total magnetic field $\vec B$ along the perimeter of the given circle is zero.}
\opt{B}{Consider two points A(3,0,0) and B(2,1,0) on the given circle. The path integral $\int_A^B \vec B\cdot d\vec\ell$ of the total magnetic field $\vec B$ along the perimeter of the given circle from A to B is $\dfrac{\mu_0 I}{2\pi}\tan^{-1}\left(\dfrac12\right)$.}
\opt{C}{The maximum value of the path integral $\int \vec B\,d\vec\ell$ of the total magnetic field $\vec B$ along the perimeter of the given circle between any two points on the circle is $\dfrac{\mu_0 I}{6}$.}
\opt{D}{The maximum value of the path integral $\int \vec B\,d\vec\ell$ of the total magnetic field $\vec B$ along the perimeter of the given circle between any two points on the circle is $\dfrac{\mu_0 I}{3}$.}
\end{qoptions}

\qn{5} In a photoelectric experiment arrangement light of frequency $\nu$ is incident on a metal target whose work function is $\phi = \dfrac{h\nu}{3}$ as shown.
\qfig{p2_phy_s1_q05_1}
\begin{qoptions}
\opt{A}{Maximum kinetic energy of photoelectrons just after emission from target is $\dfrac{2h\nu}{3}$.}
\opt{B}{Kinetic energy of photoelectrons just after emission from target may have any value ranging from zero to $\dfrac{2h\nu}{3}$}
\opt{C}{Kinetic energy of photoelectrons when they are half way between the target and collector ranges from zero to $\dfrac{2h\nu}{3}$.}
\opt{D}{Kinetic energy of photoelectrons as they reach to the collector is zero.}
\end{qoptions}

\qn{6} The energy, the magnitude of linear momentum and orbital radius of an electron in a hydrogen atom corresponding to the quantum number E, p and r respectively. Then according to Bohr's theory of hydrogen atom.
\begin{qoptions}
\opt{A}{Epr is proportional to $\dfrac1n$}
\opt{B}{$\dfrac{p}{E}$ is proportional to n}
\opt{C}{Er is constant for all orbitals}
\opt{D}{pr is proportional to n.}
\end{qoptions}

\sect{Section-II: Numerical Value}

\qn{1} There are two sets of parallel plates, one of the shape of equilateral triangle and other of pentagon (see fig.). The plates carry uniformly distributed unlike charges of equal magnitude (on each set). Deep inside the region between the plates, electric field is uniform but near the edges it becomes non--uniform. Find ratio of electric field at the mid point of the line AB in pentagon set to equilateral triangle set.
\qfig{p2_phy_s2_q01_1}

\qn{2} Five resistors are connected to a current source. Three of them have a resistance of $1\Omega$, the other two have a same but an unknown resistance. The voltage of the current source is $U_0 = 14$ V and its current is $I_0 = 10$ A. The voltage and the current strength on the third resistor are respectively $U_3 = 2$ V and $I_3 = 2$ A. The direction of the current in the resistor $R_3$ is shown in the drawing. The unknown resistance (in ohm) is:
\qfig{p2_phy_s2_q02_1}

\qn{3} Three right angled prism are glued as shown in the figure. An incident ray passes undeviated through the prism. Determine the refractive index of middle prism if refractive indices of left and right prism are $\sqrt2$ and $\sqrt3$ respectively.
\qfig{p2_phy_s2_q03_1}

\qn{4} Two identical beams A and B of plane coherent waves of the same intensity and wavelength 700 nm fall on a plane screen. The direction of beam propagation make angles $\theta_1$ and $\theta_2$ with normal to the screen and lie in same plane as shown in fig. Find distance between adjacent interference fringe on screen in $\mu$m (Here $\theta_1 = 30^\circ$ and $\theta_2 = 60^\circ$).
\qfig{p2_phy_s2_q04_1}

\qn{5} The relative density of a material is found by weighing the body first in air and then in water. The weight in air is found to be $(10.0 \pm 0.1)$ N and weight in water is found to be $(5.0\pm 0.1)$ N. From this find out maximum percentage error in relative density.

\qn{6} Determine the resistance $R_{AB}$ between A and B formed by nine identical wires of resistance $1\Omega$ each.
\qfig{p2_phy_s2_q06_1}

\sect{Section-III: Single Digit Integer (0-9)}

\qn{1} Spool of a bobbin rolls without slipping on a horizontal track. The radius of the spool is 50 cm and that of flanges is significantly more than radius of spool. At some instant speeds of two diametrically opposite points on the spool are $\sqrt7$ m/s and 3 m/s. Find out speed (in m/s) of the centre of the bobbin

\qn{2} Two particles A and B repel each other by a force $F = 8.0$ N when distance between them is less than $r_0 = 2.0$ m and apply no force on each other when distance between them is greater than $r_0$. Mass of particle A is $m = 2.0$ kg and mass of particle B is $8.0$ kg. If the particle B is projected with velocity $v = 4$ m/s towards particle A kept at rest in free space and if minimum distance between them in subsequent motion is $r_{min}$, then fill $5r_{min}$.

\qn{3} Two identical frictionless slide A and B each of height 0.5 m and mass $M = 2$ kg are placed on a horizontal frictionless floor. A disc C of mass $m = 4$ gm is released on the top of the slide A. If speed acquired by slides a long time after is V m/s, Fill 40 V.
\qfig{p2_phy_s3_q03_1}

\qn{4} A house having dimension $a = 2$m, $b = 5$m and $c = 3$m (Cross sectional view fig.).
\qfig{p2_phy_s3_q04_1}
If the least possible speed to topple over the roof thrown from anywhere left from the ground is V m/s, Find $\dfrac{V}{2}$ m/s (use $g = 10$ m/s$^2$).

\qn{5} A ball of mass 8 kg is suspended from a block of mass 12 kg with light inextensible cord. The block can slide on a frictionless slope of inclination $\theta = 30^\circ$. Initially the bar is held motionless so that the ball also stays motionless as shown in the figure. When the system is released, if initial acceleration of ball relative to ground is found to be $\dfrac{25}{n}$ m/s$^2$. Find n. (Use $g = 10$ m/s$^2$)
\qfig{p2_phy_s3_q05_1}

\qn{6} An ideal monoatomic gas is confined in a cylinder with the help of a piston that can slide inside the cylinder without friction. Walls of the cylinder have finite conductivity. Initially the piston is in equilibrium with its external surface exposed to the atmosphere at pressure $10^5$ Pa. A heat is slowly injected into the cylinder through its walls causing the gas to expand until its volume changes by 0.4 liter. If H is the amount of heat supplied (in J) find $\left(\dfrac{H}{20}\right)$.

\partsec{2}{CHEMISTRY}
\sect{Section-I: One or More Correct}

\qn{1} For the given reaction sequence.
\qfig{p2_chem_s1_q01_1}
Find out the correct option:
\begin{qoptions}
\opt{A}{Reaction involve Carbocation as an intermediate (during X--A conversion)}
\opt{B}{Reaction involve Aromatic electrophilic substitution reaction}
\opt{C}{Final product 'B' is an Ketone}
\opt{D}{The above reaction is Gattermann Aldehyde synthesis}
\end{qoptions}

\qn{2} In the given reaction the intermediate which is not involved is/are
\qfig{p2_chem_s1_q02_1}
\begin{qoptions}
\opt{A}{\qfigi{p2_chem_s1_q02_A}}
\opt{B}{\qfigi{p2_chem_s1_q02_B}}
\opt{C}{\qfigi{p2_chem_s1_q02_C}}
\opt{D}{\qfigi{p2_chem_s1_q02_D}}
\end{qoptions}

\qn{3} Which of the following are correct
\begin{qoptions}
\opt{A}{The brown ring complex obtained in the test of $\mathrm{NO_3^-}$ ion has 5 electrons in $t_{2g}$}
\opt{B}{In $[\mathrm{Fe(NH_3)_6}]^{+2}$, $P > \Delta_0$ hence all electrons remain in $e_g$.}
\opt{C}{For $[\mathrm{NiF_6}]^{-2}$, $\Delta_0 > P$ hence it's diamagnetic by nature}
\opt{D}{In $[\mathrm{PtCl_3(C_2H_4)}]^-$, one Pt--Cl bond is stronger than remaining two}
\end{qoptions}

\qn{4} During Hydrolysis of SiF$_4$, which are not true?
\begin{qoptions}
\opt{A}{It undergoes complete hydrolysis}
\opt{B}{On per mole of SiF$_4$ consumed, 50\% molecules are hydrolysed and rest 50\% produce hexaflurosilicic acid}
\opt{C}{Out of total 3 moles SiF$_4$ consumded, only two moles gets hydrolysed}
\opt{D}{Among silicon halides, SiF$_4$ has maximum ease of hydrolysis}
\end{qoptions}

\qn{5} Which of the following are correct statements?
\begin{qoptions}
\opt{A}{Mond's process \& Van--Arkel de Boer's process are examples of vapour phase refining}
\opt{B}{In Parke's process of desilverisation of lead, zinc remains in lead and Ag being more volatile is removed.}
\opt{C}{During electrolytic refining of Cu; the anode mud contains Ag, Au, Pt, Se, Te, Sb etc. in oxidised form}
\opt{D}{During polling process of Sn extraction, impurity of SnO$_2$ present in metal is reduced to metallic tin.}
\end{qoptions}

\qn{6} Which of the following is/are correct?
\begin{qoptions}
\opt{A}{Langmuir adsorption theory states that reactivity of a site also depends on the vacancy status of nearby sites.}
\opt{B}{Order of adsorption can be changed by changing pressure conditions.}
\opt{C}{Desorption of gas on solid is favoured by low temperature.}
\opt{D}{ZSM--5 is shape independent catalyst used to crack hydrocarbons.}
\end{qoptions}

\sect{Section-II: Numerical Value}

\qn{1} (i) $\mathrm{HC \equiv CH + CH_2N_2 \xrightarrow{dark} A}$ (Degree of unsaturation is 'X').\\
(ii) $\mathrm{CH_2=CH_2 + CH_2N_2 \xrightarrow{\Delta} B}$ (Degree of unsaturation is 'Y')\\
(iii)
\qfig{p2_chem_s2_q01_1}
$\longrightarrow$ C. (Degree of unsaturation is 'Z').\\
Find out value of $\dfrac{Z(X)+Y}{X}$.

\qn{2} Galactose like glucose exist in two isomeric cyclic form $\alpha$--D Galactopyranose 28\% and $\beta$--D Galactopyranose 72\%. $\beta$--D Galactopyranose being more stable isomer and show specific rotation of $+52.8^\circ$ where as $\alpha$--D Galactopyranose shows specific rotation of $150.7^\circ$. Find out the sp. rotation of mixture at equilibrium.

\qn{3} What limiting amount of Na$_2$CO$_3$ (in millimoles) is to be added to 150 ml of 0.1 M carbonated water to make its pH = 6.4? Given: $K_{a_1}$ and $K_{a_2}$ for H$_2$CO$_3$ = $4\times10^{-8}$ \& $4\times10^{-12}$ respectively.

\qn{4}
\qfig{p2_chem_s2_q04_1}
$\xrightarrow[\text{(iv) Electrolysis}]{\substack{\text{(i) Excess KCN}\\\text{(ii) H}_3\text{O}^+\\\text{(iii) Aq. NaOH}}}$ Product\\
Total no. of organic products = x\\
Total Faradays of electricity required to electrolyse 2.44 gm product obtained in step (iii) assuming 15\% current efficiency = y\\
Find $(x+y)$

\qn{5} No. of $e^-$ in ABMO of cationic part in O$_2$SbF$_6$ = x\\
No. of nodal planes in $\delta$--Bond ABMO formation = y\\
No. of oxygen atoms per tetrahedral unit in amphibole silicate = z\\
No. of hybridised orbitals in central atom of ClO$_2$ = w\\
find $(x+y+z-w)$

\qn{6} $[\mathrm{M(gly)a_2b_2}] \xrightarrow{\text{one of the ligand b is subtituted by ligand c}} \text{product complex.}$\\
Find sum of geometrical isomers of reactant complex and product complex.

\qfullbegin
\sect{Section-III: Single Digit Integer (0-9)}

\qn{1} How many reaction given below are correctly matched with their products.
\qfig{p2_chem_s3_q01_1}
\qn{2} How many of the following reactions can give Benzene as final product.
\qfig{p2_chem_s3_q02_1}
\qfullend

\qn{3} No. of correct statements from the following = x. Find x.\\
(i) K$_2$CO$_3$ is prepared by Solvey's process\\
(ii) Na$_2$CO$_3$ on reaction with MgCl$_2$ gives basic magnesium carbonate.\\
(iii) Na$_2$CO$_3\cdot$10H$_2$O is effloresces.\\
(iv) F$_2$ on reaction with NaOH gives OF$_2$\\
(v) NaOH on reaction with HgCl$_2$ gives yellow precipitate.\\
(vi) Elements like B,Si disproportionate with NaOH\\
(vii) NaOH has higher thermal stability as compared to KOH.\\
(viii) Dilute solution of Na in liq. NH$_3$ has lower density than pure solvent\\
(ix) In castner--Kellner's cell, hydrogen has low overpotential on Hg due to which Na is deposited on Hg cathode.

\qn{4} What will be the total surface area (cm$^2$) of a regular truncated octahedron of edge length 0.5 cm? (Take: $4\sqrt3 = 7.0$)

\qn{5} 100 ml CuSO$_4$ (aq.) solution was titrated for 1250 s, using a current of 0.965A with the same efficiency which produces 8.64 gm Ag [Ag = 108] from AgNO$_3$ using 10A current for 965 s. Find final pH of CuSO$_4$ solution.

\qn{6} Enthalpy of formation, $\Delta H_f(\mathrm{Fe_3O_4}) = -1092$ kJ/mol.\\
Enthalpy of sublimation, $\Delta H_{sub}(\mathrm{Fe}) = 390$ kJ/mol.\\
B.D.E. of O$_2$ = 490 kJ/mol.\\
I.E.$_1$ of Fe = 760 kJ/mol.\\
I.E.$_2$ of Fe = 1560 kJ/mol.\\
I.E.$_3$ of Fe = 2960 kJ/mol.\\
I.E.$_1$ of O$^{2-}$ = $-844$ kJ/mol.\\
I.E.$_2$ of O$^{2-}$ = 142 kJ/mol.\\
Magnitude of Lattice Enthalpy of Fe$_3$O$_4$ = 'x' kJ/mol. Find the value of $\dfrac{x}{3155}$.

\partsec{3}{MATHEMATICS}
\sect{Section-I: One or More Correct}

\qn{1} Consider tetrahedron whose three coterminous edges are represented by $\vec a\times\vec b$, $\vec b\times\vec c$ and $\vec c\times\vec a$, where
$$\vec a = (2+\sin\theta)\hat i + (\cos\theta)\hat j + (\sin2\theta)\hat k$$
$$\vec b = \sin\left(\theta+\dfrac{2\pi}{3}\right)\hat i + \cos\left(\theta+\dfrac{2\pi}{3}\right)\hat j + \sin\left(2\theta+\dfrac{4\pi}{3}\right)\hat k$$
$$\vec c = \sin\left(\theta-\dfrac{2\pi}{3}\right)\hat i + \cos\left(\theta-\dfrac{2\pi}{3}\right)\hat j + \sin\left(2\theta-\dfrac{4\pi}{3}\right)\hat k$$
Where $\theta\in\left(0,\dfrac{\pi}{2}\right)$. Let $f(\theta)$ represent volume of tetrahedron, then
\begin{qoptions}
\opt{A}{No. of solution of $f(\theta) = \dfrac14$ is 3}
\opt{B}{Maximum value of $f(\theta)$ is 2}
\opt{C}{$\displaystyle\sum_{n=1}^{5} f\left(\dfrac{n\pi}{33}\right) = \dfrac{9}{8}$}
\opt{D}{No. of non-differentiable points of $\sqrt{f(\theta)}$ is 1}
\end{qoptions}

\qn{2} The solution of $\left(\dfrac{dy}{dx}\right)^2 - 2\left(\dfrac{x^2+y^2}{x^2-y^2}\right)\dfrac{dy}{dx} + 1 = 0$ may be where C is any arbitrary constant
\begin{qoptions}
\opt{A}{$y^2+2xy-x^2 = C$}
\opt{B}{$2\tan^{-1}\left(\dfrac{y}{x}\right) = \ln(x^2+y^2) + C$}
\opt{C}{$2\tan^{-1}\left(\dfrac{x}{y}\right) = \ln(x^2+y^2) + C$}
\opt{D}{$x^2+2xy-y^2 = C$}
\end{qoptions}

\qn{3} Let PQ be a focal chord of a parabola. A circle with diameter PQ always touches the line $x+5=0$ for all positions of P and Q. Length of latus rectum of parabola is 8. If equation of line PQ is $x+y=1$, then
\begin{qoptions}
\opt{A}{Parabola is $(y-2)^2 = 8(x+3)$}
\opt{B}{Parabola is $(y-10)^2 = -8(x+7)$}
\opt{C}{From point $(2,3)$ two $\perp$ normal can be drawn to one of the possible parabola}
\opt{D}{From the point $(-17,12)$ two $\perp$ normal can be drawn to one of the possible parabola}
\end{qoptions}

\qn{4} Let A, B and P be three non-singular matrices of order 2 satisfying $AB - BP = 2I$ and $\det(P) = 3$, $\mathrm{Tr}(P) = 4$ and $B = \mathrm{adj}(B)$ where I = Identity matrix, $\det(P)$ = determinant of matrix P, $\mathrm{Tr}(P)$ = Trace of matrix P
\begin{qoptions}
\opt{A}{If $\det(A) = 15$, then $\mathrm{Tr}(B) = 2$}
\opt{B}{If $\det(A) = 15$, then $\mathrm{Tr}(B) = -2/3$}
\opt{C}{If $\mathrm{Tr}(B) = 4$, then $\det(A) = 4$}
\opt{D}{If $\mathrm{Tr}(B) = 4$, then $\mathrm{adj}(\mathrm{adj}(A)) = 8$}
\end{qoptions}

\qn{5} Let $f:\mathbb{R}\to\mathbb{R}$, $g:\mathbb{R}\to\mathbb{R}$ and $h:\mathbb{R}\to\mathbb{R}$ be differentiable function such that $f(x) = x^3+5x+1$, $g(f(x)) = x$ and $h(g(g(x))) = x\ \forall x\in\mathbb{R}$ then
\begin{qoptions}
\opt{A}{$h(0) = 7$}
\opt{B}{$h'(1) = 1216$}
\opt{C}{$g'(1) = \dfrac15$}
\opt{D}{$h(g(2)) = 19$}
\end{qoptions}

\qn{6} There are 20 points lying on a circle, distance between two consecutive points is same. Three points are selected at random. X represent no. of ways to select them. Y represent no. of ways to select them such that no two adjacent objects are selected and Z represent no. of ways to select such that chosen objects are not diametrically opposite then
\begin{qoptions}
\opt{A}{X = 1140}
\opt{B}{Y = 800}
\opt{C}{Z = 960}
\opt{D}{$Y\cap Z = 660$}
\end{qoptions}

\sect{Section-II: Numerical Value}

\qn{1} Let $100n\,a_n = 100(n-2)a_{n-1} + 1$ and $a_0 = 0$ for $n\ge1$ then $\left[\displaystyle\sum_{n=0}^{4044} a_n\right]$, where $[\ ]$ denotes greatest integer less than or equal to x, is equal to?

\qn{2} A function of $f(x_1,x_2,x_3,\ldots,x_n)$ is said to be linear in each of its variables if it is a polynomial such that no variables appears with power higher than one in any term. For example $1+x+y+2xy$ is linear in x and y but $1+y^2$ is not. A function $f(x,y)$ is linear in x and y and $f(x,y) = \dfrac{9}{xy}\ \forall x,y\in\{3,4\}$ then $f(5,5)$ is

\qn{3} If\\
\resizebox{\linewidth}{!}{$I = \displaystyle\int_1^3\left(\tan^{-1}\left(\dfrac{x^2-5x+6}{x^3-6x^2+12x-7}\right) + \tan^{-1}\left(\dfrac{1}{x^2-4x+4}\right)\right)dx$}\\
then I is

\qn{4} There are 3 different pair (i.e. 6 unit say $a_Ra_L$, $b_Rb_L$, $c_Rc_L$) of shoes in a lot. Now three person came and pick the shoes randomly (each gets a 2 unit), then the probability that no one is able to wear shoes (i.e. no one gets correct pair)

\qn{5} If the sum of the coefficient of all even power of x in the product $(1+x+x^2+\ldots+x^{2n})(1-x+x^2-x^3+\ldots+x^{2n})$ is 81 then n is

\qn{6} Let p and q be the length of two chords of a circle which subtend angles $36^\circ$ and $60^\circ$ respectively at the centre of the circle and the angle subtended by the chord of length $(p+q)$ at the center of circle is $k^\circ$ then k is

\sect{Section-III: Single Digit Integer (0-9)}

\qn{1} In a triangle ABC, the median from B and C are perpendicular, then minimum integral value of $(\cot B + \cot C)$ is

\qn{2} Let $\dfrac{1}{K_r - 2i}$ where $K_r\in\mathbb{R}$ for $r = 1,2,3,4,5,6,7,8$ are vertices of a octagon in Argand plane. If area of circle inscribed in this octagon is $\dfrac{\pi(\sqrt a+1)}{b\sqrt2}$, then $\dfrac{b}{2a}$ is (where $i=\sqrt{-1}$)

\qn{3} Consider $f(x) = \dfrac{x^2+4}{|x|}$ and $g(x) = p|x|+2$, $p>0$. If the exhaustive set of values of 'p' for which $f(x) > g(x)\ \forall x\in\mathbb{R}-\{0\}$ is $(a,b)$ then $a+4b$ is

\qn{4} Consider the ellipse $\dfrac{x^2}{25} + \dfrac{y^2}{16} = 1$. Let A, B, C, D be four points on this ellipse such that the normal drawn at these points are concurrent at $(3,6)$. The center of the conic other than given ellipse on which these four points lies is $(p,q)$ then $\left(\dfrac{p-q}{19}\right)$ is equal to

\qn{5} A perpendicular is drawn from a point on the line $\dfrac{x-1}{2} = \dfrac{y+1}{-1} = \dfrac{z}{1}$ to the plane $x+y+z=3$ such that the foot of the perpendicular Q also lies on the plane $x-y+z=3$, if the coordinate of Q is $(a,b,c)$ then $|a-b-c|$ is

\qn{6} If $\theta$ is the angle between $\vec a$ and $\vec b$ such that $\sqrt3(\vec a\cdot\vec b) + \left|\vec b - (\vec a\cdot\vec b)\vec a\right| = 1 + \vec b\cdot\vec b$ where $|\vec a| = 1$ then $\mathrm{cosec}^2\theta$ is

\end{multicols}
\end{document}
